\section{Dynamical inputs and standard models}
\label{sec:inputs}

Fix an integer $n\ge2$ and put $I=\{1,\ldots,n\}$. We first state
the dynamical results used below. The distinction between norm and
pointwise convergence matters: norm convergence is available for any
finite number of slots, whereas the pointwise input on arbitrary
systems will be used only for at most three slots.

\subsection{Standard models and the two companion inputs}

For the fixed tuple in Theorem~\ref{thm:main}, record all integer
translates of bounded measurable representatives of the functions.
This gives a measurable map into a countable product of compact scalar
discs, intertwining $T$ with the invertible shift. The pushforward
measure defines a standard Borel probability system, and mixing passes
to this factor. The original functions are its coordinate functions.
Changing representatives affects at most a countable union of translated
null sets, so almost-everywhere convergence on the factor pulls back.
We may therefore work on standard probability spaces.

Throughout the proof, factors and conditional laws are understood
modulo null sets. We use Borel versions of measurable functions.
Sub-$\sigma$-algebras on a standard probability space have standard
models modulo null sets. A flow is a jointly Borel action
$(S_Y^t)_{t\in\R}$ on a standard Borel probability space $(Y,m_Y)$,
preserving $m_Y$. Factor maps commute with every fixed time modulo
null sets. Joint measurability permits Fubini in all time integrals.
For a preserved factor, induced transformations at fixed times can be
modeled as invertible maps; we take time integrals on the original flow.

We use the following statement from
\cite[Theorem~1.1]{MixingCompanion}.
\begin{theorem}[Higher-order mixing, companion input]
\label{thm:multiple-mixing}
Let $T$ be an invertible mixing probability-preserving transformation.
For any fixed $s\ge2$, $g_1,\ldots,g_s\in L^\infty(\mu)$,
and integer times $t_1,\ldots,t_s$,
\[
 \int_X\prod_{j=1}^s g_j(T^{t_j}x)\dd\mu(x)
 \longrightarrow\prod_{j=1}^s\int_Xg_j\dd\mu
 \quad\text{as}\quad
 \min_{i\ne j}|t_i-t_j|\longrightarrow\infty.
\]
For $s=1$ the equality holds for every time.
\end{theorem}
The cited theorem is stated for sets and ordered times. Simple-function
approximation, invariance, and the finitely many orderings of the times
give this formulation. No quantitative mixing rate is asserted or used.

Our second companion input is the smooth-average convergence consequence
of \cite[Proposition~3.5 and Lemma~3.6]{SlopesCompanion}.
An \emph{averaging profile} there is a real compactly supported smooth
function $\varphi$ of mass one with moments of orders $1,2,3$ equal to
zero, and, for some $G\ge1$,
\[
 \norm{\varphi^{(r)}}_\infty
 \le \bigl(G(r+2)\bigr)^{G(r+2)}\qquad(r\ge0).
\]
The constant $G$ may depend on the profile. We need the following
two precise conclusions.
\begin{proposition}[Smooth triple averages, companion input]
\label{prop:smooth-input}
Let $R$ be any invertible probability-preserving transformation, and
let $b_1,b_2,b_3$ be distinct nonzero integers with greatest common
divisor one. For bounded $g_1,g_2,g_3$ and every averaging profile,
\[
 V_k^\varphi(y)=\sum_{l\in\Z}2^{-k}\varphi(l/2^k)
              \prod_{j=1}^3g_j(R^{b_jl}y)
\]
converges almost everywhere as $k\to\infty$. Moreover, for each
$c\in[1,2]$ and $\epsilon>0$ there is such a profile with
\[
 \norm{\varphi-c^{-1}\mathbf1_{(0,c]}}_{L^1(\R)}<\epsilon.
\]
\end{proposition}
The approximation is \cite[Lemma~3.9]{SlopesCompanion}. Neither part
requires mixing or ergodicity. We will not assume pointwise convergence
of longer averages on arbitrary systems.

\subsection{Relative weak mixing and distal factors}

We recall the structure facts in the form needed for the auxiliary
systems. An extension of an invertible system $R$ over a factor $B$
is \emph{relatively weakly mixing} if, for bounded $f,g$ with
$\E[f\mid B]=0$,
\begin{equation}\label{eq:rwm}
 \frac1H\sum_{h=1}^H
 \norm{\E[\overline f(g\circ R^h)\mid B]}_2^2\longrightarrow0.
\end{equation}
An extension is \emph{relatively compact} if the invariant finitely
generated $L^\infty(B)$-modules span a dense subspace of its $L^2$
space. A module consists of finite sums $\sum_j a_jF_j$, with
$a_j\in L^\infty(B)$ and fixed $F_j\in L^2$; invariance is under
composition with both $R$ and $R^{-1}$. Individual modules need not be
closed; see the corrected module characterization in
\cite[Theorem~2.5(ii)$'$]{JamneshanCorrigendum2026}.
A \emph{relatively distal} extension is a tower of relatively
compact extensions and inverse limits. A system is \emph{distal} if
it is distal over the trivial factor. A joining is an invariant
probability law on a product with the prescribed marginals.

\begin{proposition}[Relative structure]\label{prop:structure}
The following statements hold for standard probability-preserving
systems, without an ergodicity assumption.
\begin{enumerate}
\item Every system $Y$ has a maximal measure-distal factor $W_Y$,
and $Y\to W_Y$ is relatively weakly mixing. Every distal factor
pulls into the maximal distal factor of an extension. Automorphisms
commuting with the transformation preserve $W_Y$.
\item Factors and joinings of distal systems are distal. A distal
system is relatively distal over each of its factors. Distality is
equivalent for a transformation and any of its nonzero integer powers.
\item Relative weak mixing passes to nonzero integer powers. The
conditional product of relatively weakly mixing fibre laws over an
invariant joining of their base laws is relatively weakly mixing
over that joining.
\item A relatively weakly mixing extension and a relatively distal
extension of the same base are relatively disjoint: every invariant
joining identifying their bases is the conditional product over
that base.
\item For a distal transformation $R$, every $d\ge1$, and bounded
$g_1,\ldots,g_d$, the averages
\[
 \frac1N\sum_{l=1}^N\prod_{i=1}^d g_i\circ R^{il}
\]
converge almost everywhere.
\end{enumerate}
For a flow, the maximal distal factor of time one is preserved by
the flow, and its action at every nonzero rational time is distal.
\end{proposition}

These are the conclusions of
\cite[Proposition~2.2]{FourfoldCompanion}. They include the
nonergodic form of the Furstenberg--Zimmer structure theorem
\cite{Zimmer1976,Jamneshan2023} and the pointwise theorem for
distal systems \cite[Theorem~C]{HuangShaoYe2019}. The cited proposition
gives the component-disintegration argument needed for the latter
in the nonergodic setting. In the flow assertion, invariance of
$W_Y$ follows because each time commutes with time one. If
$r=p/q\ne0$ is rational, the $q$th power of its action is the
distal action at integer time $p$, so the power assertion applies.

The following short argument, from
\cite[Lemma~2.3]{FourfoldCompanion}, will also fix the norm-convergence
input explicitly.
\begin{lemma}[Relative norm convergence]\label{lem:relative-norm}
Suppose $Y\to B$ is relatively weakly mixing for $R$. For distinct
nonzero integers $a_1,\ldots,a_d$ and bounded $G_1,\ldots,G_d$,
\begin{equation}\label{eq:relative-norm}
 \left\|\frac1N\sum_{l=1}^N
 \left(\prod_{i=1}^dG_i\circ R^{a_il}
       -\prod_{i=1}^d\E[G_i\mid B]\circ R^{a_il}\right)\right\|_2
 \longrightarrow0.
\end{equation}
\end{lemma}
\begin{proof}
By telescoping and scaling, it suffices to show that the average of
$v_l=\prod_i G_i\circ R^{a_il}$ tends to zero when all functions
are bounded by one and one of them, $G_s$, is centered over $B$.
Induct on $d$. For fixed $h\ge1$ put
$F_{i,h}=\overline{G_i}(G_i\circ R^{a_ih})$. Invariance gives
\[
 \ip{v_l}{v_{l+h}}
 =\int_Y F_{s,h}\prod_{i\ne s}
                  F_{i,h}\circ R^{(a_i-a_s)l}\dd m_Y.
\]
The induction hypothesis replaces the factors indexed by $i\ne s$
in their average by their conditional expectations over $B$. The
remaining product is $B$-measurable and bounded by one; hence
\[
 \limsup_{N\to\infty}
 \left|\frac1N\sum_{l=1}^N\ip{v_l}{v_{l+h}}\right|
 \le\norm{\E[F_{s,h}\mid B]}_1.
\]
For $d=1$ this bound holds directly. The Hilbert-space van der
Corput inequality now bounds the limiting squared norm by a constant
times
\[
 \frac1H+\frac1H\sum_{h=1}^H\norm{\E[F_{s,h}\mid B]}_1.
\]
Relative weak mixing for $R^{a_s}$ and Cauchy--Schwarz make this
expression tend to zero as $H\to\infty$.
\end{proof}

For positive slopes, inserting functions $1$ in omitted positions
and using Proposition~\ref{prop:structure}(5) shows that the
distal-factor averages have an $L^2$ limit. Thus
Lemma~\ref{lem:relative-norm} gives an $L^2$ limit for arbitrary
bounded tuples with distinct positive slopes on every standard
system. On the original mixing system, the trivial factor already
satisfies \eqref{eq:rwm}; the limit is the product of the means.

\subsection{Triple convergence along all lengths}

\begin{lemma}[Triple pointwise convergence]\label{lem:triple-convergence}
On any standard invertible probability-preserving system, averages
of at most three bounded functions at distinct positive integer
slopes converge almost everywhere along all positive integer lengths.
\end{lemma}
\begin{proof}
Add constant-one functions if necessary to have three distinct
positive slopes, write them as $gb_1,gb_2,gb_3$ with
$\gcd(b_1,b_2,b_3)=1$, and apply
Proposition~\ref{prop:smooth-input} to $S=R^g$. By scaling, suppose
the three functions are bounded by one, and put
$F_l=\prod_{j=1}^3g_j\circ S^{b_jl}$.

For each rational $c\in[1,2]$, choose a sequence of profiles tending
in $L^1$ to $h_c=c^{-1}\mathbf1_{(0,c]}$. On a common conull set
all the corresponding smooth dyadic averages converge. For
$N_k=\lfloor c2^k\rfloor$, the difference between the ordinary
average at $N_k$ and $V_k^\varphi$ is at most
\[
 \frac1{c2^k}
 +2^{-k}\sum_{l\in\Z}|h_c(l/2^k)-\varphi(l/2^k)|.
\]
The last sum tends to $\norm{h_c-\varphi}_1$ by Riemann sums.
Arbitrarily close profiles therefore give the Cauchy property on
each of these scales. The $L^2$ limit established above identifies
all their almost-everywhere limits with the same function.

For any bounded sequence $|F_l|\le1$, changing the averaging length
from $N$ to $M$ changes the average by at most
\begin{equation}\label{eq:length-change}
                 \frac{2|M-N|}{\max(M,N)}.
\end{equation}
For $2^k\le N<2^{k+1}$, compare $N$ with
$\lfloor c2^k\rfloor$ on a finite rational grid in $[1,2]$.
Equation~\eqref{eq:length-change} makes the limiting error
arbitrarily small as the grid mesh tends to zero. This proves
convergence along all integer lengths.
\end{proof}
